% ============================================================
%  DBMS_10 – Project Proposal for the Term Project
%  THGA Bochum · Databases · Stephan Bökelmann
% ============================================================
\documentclass[11pt,a4paper]{article}

\usepackage{thga-db}

% ----- Metadata --------------------------------------------
\renewcommand{\thgaKurs}{Databases}
\renewcommand{\thgaDatum}{Summer Term 2026}

\begin{document}
\thispagestyle{empty}
\thgaTitel{Exercise 10 – Project Proposal}%
{Term project: design your own database-backed system}

\begin{infobox}
\textbf{Module:} Introduction to Database Management Systems · THGA Bochum \\
\textbf{Lecturer:} Stephan Bökelmann · \href{mailto:sboekelmann@ep1.rub.de}{sboekelmann@ep1.rub.de} \\
\textbf{Repository:} \url{https://github.com/MaxClerkwell/DBMS_10} \\
\textbf{Prerequisites:} Lectures 01--10, exercises DBMS\_01--DBMS\_09 \\
\textbf{Type of exercise:} design and written submission -- \emph{not} a coding assignment \\
\textbf{Size (proposal):} 2--4 pages \\
\textbf{Effort (term project):} max.\ 2 months window, approx.\ 40\,h of real work incl.\ documentation and video
\end{infobox}

\tableofcontents
\vspace{1em}

% ============================================================
\section{What this exercise is about}
% ============================================================

This exercise is deliberately \textbf{structured differently} from the previous
ones. This time you do not write code and you do not work through a
step-by-step tutorial. Instead you \textbf{design} your own project and submit a
\textbf{project proposal} to the lecturer.

The proposal is the basis of your \textbf{term project} -- the graded
deliverable in which you build your own database-backed system over the course
of the term. Before you start building, we settle the \emph{subject}: which
application do you want to implement? Which data does it manage? Which functions
does it offer?

\begin{infobox}
\textbf{Goal of this exercise:} Submit a 2--4-page project proposal that states
\emph{which} database-backed system you want to build for the term project and
\emph{how} it draws on the contents of lectures 01--10.
\end{infobox}

The proposal is \textbf{reviewed and approved} by the lecturer. Only after
approval do you start the implementation. That way we make sure -- before the
real effort begins -- that your undertaking has the right size (neither too
small nor too large) and covers all mandatory building blocks.

% ============================================================
\section{The term project at a glance}
% ============================================================

For the term project you build a \textbf{complete, database-backed system for an
application of your choice}. The topic is free: a club membership manager, a
recipe and shopping planner, a library, a ticketing system, a warehouse for a
model-building club, a training app -- anything is allowed as long as
meaningful data and relationships can be modelled.

Your application should reflect the chain we have built up over the term:
\emph{from data modelling to a running, usable application.}

\subsection{The mandatory technology stack}

The system must consist of the following layers. You already know each layer
from the lectures and the previous exercises:

\begin{itemize}[leftmargin=1.4em]
  \item \textbf{PostgreSQL} as the relational database management system -- your
        cleanly modelled schema with keys and integrity constraints
        (L 03--07). Runs as a \textbf{container on the lecture server}.
  \item \textbf{FastAPI} as the REST layer in Python -- the database is
        \emph{never} accessed directly by the frontend, always through the API
        (L 08). The API is protected with an \textbf{\texttt{X-API-Key}}
        (cf.\ DBMS\_09) and also runs as a \textbf{container on the lecture
        server}.
  \item \textbf{Docker Compose} orchestrates \textbf{only the backend} --
        database and API run as services that start together with
        \texttt{docker compose up} (L 09). This Compose setup is
        \textbf{deployed on the lecture server}.
  \item \textbf{A frontend} that uses the API with the \texttt{X-API-Key}. It is
        \textbf{not containerised}. Instead you build it into a \textbf{Debian
        installer (\texttt{.deb})} and install it on the lecture server, so the
        lecturer can start and test your application there directly (L 10/13). A
        \textbf{Python frontend} (e.g.\ tkinter as in L 10) is recommended; a
        different frontend is allowed but must likewise install as a
        \texttt{.deb}.
\end{itemize}

\begin{beispielbox}
\textbf{The architecture -- everything on the lecture server:}\\[6pt]
\centerline{%
  \textbf{Frontend} (installed \texttt{.deb})
  $\;\xrightarrow{\;\text{HTTP + X-API-Key}\;}\;$
  \textbf{FastAPI}
  $\;\longrightarrow\;$ \textbf{PostgreSQL}%
}\\[6pt]
The \textbf{backend} (API + database) runs in containers via \textbf{Docker
Compose}. The \textbf{frontend} runs \emph{next to it} as an installed Debian
package and talks to the API over HTTP using the \texttt{X-API-Key}.
\end{beispielbox}

\subsection{Deployment on the lecture server}

The finished system runs on the \textbf{lecture server} so the lecturer can test
it without any setup of their own:

\begin{itemize}[leftmargin=1.4em]
  \item \textbf{Backend:} \texttt{docker compose up -d} starts database and API
        as containers.
  \item \textbf{Frontend:} installed as a \texttt{.deb}
        (\texttt{sudo dpkg -i \dots}) and then launched via the application
        launcher or \texttt{/usr/bin/\dots}.
  \item \textbf{Access control:} write calls to the API require the correct
        \texttt{X-API-Key}; the frontend asks for the API URL and key at startup
        (connection dialog, L 10).
\end{itemize}

\subsection{What a sensible size is}

Plan for a data model of roughly \textbf{three to six tables} with at least one
\texttt{N:M} relationship, a handful of API endpoints (reading \emph{and}
writing) and a frontend that operates the most important of those endpoints.
That is enough to demonstrate all the concepts you have learned, and small
enough to finish within the given effort.

% ============================================================
\section{Deliverables and effort}
% ============================================================

The finished term project consists of \textbf{three} parts, submitted together:

\subsection{The running system}

Database and API as containers on the lecture server (via Docker Compose), the
frontend installed as a \texttt{.deb} on the same server. The source code lives
in a Git repository.

\subsection{The written documentation (GitHub repo with LaTeX CI)}

The documentation is itself a small \textbf{LaTeX project} that describes the
undertaking, the data model, the API, the frontend and the operation -- similar
to the handouts of this course. It explains your design decisions and shows how
the concepts from L 01--10 were implemented.

\textbf{Submission format:} a \textbf{GitHub repository} with an integrated
LaTeX build workflow (GitHub Actions) and a \texttt{Makefile}. On every tag push
the workflow builds the PDF and attaches it to a \textbf{GitHub Release}.

\begin{infobox}
\textbf{You do not have to invent this setup.} This repository (DBMS\_10) ships a
ready-to-use workflow (\texttt{.github/workflows/build.yml}) and \texttt{Makefile}
-- \textbf{copy them} into your own documentation repository. A complete,
buildable example lives in \texttt{example-documentation/}, and a fill-in
skeleton for the proposal in \texttt{proposal-template/}.
\end{infobox}

\subsection{The video (8--10 minutes)}

Record an \textbf{8--10-minute video} in which you demonstrate your system and
explain the most important design decisions. For the submission you have
\textbf{two options}:

\begin{itemize}[leftmargin=1.4em, topsep=2pt]
  \item upload it as an \texttt{.mpg} file to \textbf{Moodle}, \emph{or}
  \item upload it to \textbf{YouTube}, set the visibility to
        \textbf{``unlisted''} and submit only the \textbf{link}.
\end{itemize}

\subsection{Time frame and effort}

\begin{itemize}[leftmargin=1.4em, topsep=2pt]
  \item \textbf{Working window:} at most \textbf{2 months} from approval of the
        proposal.
  \item \textbf{Real workload:} about \textbf{one working week}
        ($5 \times 8\,\text{h} \approx 40\,\text{h}$) --
        \emph{including} documentation and video.
\end{itemize}

Scope your project to this budget: the system, the documentation \emph{and} the
video together should be finished in roughly 40 hours of work.

\subsection{Reuse of code}

You \textbf{may reuse} code examples from the \textbf{lectures and the practical
exercises}, provided you \textbf{cite the source} (e.g.\ ``after L 09, Docker
Compose section'' or ``taken from DBMS\_08''). Third-party code without a source
citation counts as an attempt to deceive.

% ============================================================
\section{What the lectures contribute}
% ============================================================

Your proposal -- and later the term project -- should visibly build on the
contents of lectures \textbf{01--10}. The table below shows which building block
comes from which lecture and how it becomes visible in the project.

\renewcommand{\arraystretch}{1.3}
\begin{tabularx}{\linewidth}{@{}p{1.2cm}p{4.4cm}X@{}}
\toprule
\textbf{L} & \textbf{Topic} & \textbf{How it becomes visible in the project} \\
\midrule
01 & What is a DBMS, three-level model & Justification \emph{why} a relational
     DBMS fits the undertaking \\
02 & ER modelling, cardinalities & An ER sketch with entities, relationships and
     cardinalities \\
03 & Relational model, keys & Primary and foreign keys, referential integrity \\
04 & Normalisation, normal forms & A schema at least in 3NF, free of redundancy \\
05 & SQL DDL/DML & \texttt{CREATE TABLE} script with constraints; seed data \\
06 & SQL queries, JOINs, aggregation & Queries behind the API endpoints (at
     least one JOIN, one aggregation) \\
07 & PostgreSQL in operation & Running on PostgreSQL, init script, data types \\
08 & Python, psycopg2, FastAPI & The REST layer with read and write endpoints,
     protected by \texttt{X-API-Key} \\
09 & Docker, Compose, .env & \texttt{docker-compose.yml} (postgres + api),
     secrets in \texttt{.env}, deployment on the lecture server \\
10 & Frontend basics & A usable frontend with a connection dialog (URL +
     \texttt{X-API-Key}), built as a \texttt{.deb} installer \\
\bottomrule
\end{tabularx}

\begin{infobox}
\textbf{What is \emph{not} mandatory:} The following are \textbf{optional},
voluntary extensions:
\begin{itemize}[leftmargin=1.4em, topsep=2pt]
  \item installers for further platforms (Windows setup \texttt{.exe}, macOS
        \texttt{.dmg}) -- the Debian installer (\texttt{.deb}) for the frontend,
        by contrast, is \textbf{required};
  \item advanced transaction and isolation behaviour (ACID, isolation levels,
        locks, MVCC) -- L 11;
  \item time-series data, InfluxDB and Grafana -- L 12.
\end{itemize}
Whoever includes these topics anyway is welcome to do so as a voluntary
extension.
\end{infobox}

% ============================================================
\section{The project proposal -- your submission}
% ============================================================

Submit a PDF of \textbf{2--4 pages} containing the following sections. Use the
headings as a template -- a ready-to-fill skeleton is provided in
\texttt{proposal-template/}.

\subsection{Application domain and motivation}

Describe in a few sentences \emph{which} application you build and \emph{which
problem} it solves. Who uses it, and for what? One or two concrete usage
scenarios make the proposal tangible.

\subsection{Data model -- ER sketch}

Draw an \textbf{ER sketch} (L 02) with the central entity types, their most
important attributes and the relationships including \textbf{cardinalities}. A
clean hand drawing photographed is entirely sufficient. Make sure at least one
\texttt{N:M} relationship occurs.

\subsection{From model to schema}

Derive a \textbf{relational schema} from the ER model (L 03). For each table,
name the columns, the \textbf{primary key} and the \textbf{foreign keys}.
Briefly argue that your schema is free of redundancy (at least 3NF, L 04) -- one
or two sentences suffice.

\subsection{API design}

Sketch the planned \textbf{FastAPI endpoints} (L 08). A table is enough:

\renewcommand{\arraystretch}{1.2}
\begin{tabularx}{\linewidth}{@{}p{2cm}p{4cm}X@{}}
\toprule
\textbf{Method} & \textbf{Path} & \textbf{Purpose} \\
\midrule
\texttt{GET}  & \texttt{/…}  & reads data (with JOIN/aggregation behind it) \\
\texttt{POST} & \texttt{/…}  & creates a new record (requires \texttt{X-API-Key}) \\
\multicolumn{3}{@{}l}{\itshape … at least one reading and one writing endpoint} \\
\bottomrule
\end{tabularx}

Mark behind which endpoint a \textbf{query with a JOIN} and where an
\textbf{aggregation} sits (L 06), and which endpoints require the
\texttt{X-API-Key}.

\subsection{Frontend idea}

Briefly describe what the frontend looks like and which endpoints it operates
(L 10). A rough wireframe or a list of the screens/forms is enough. State
whether you use the recommended Python variant (tkinter) or something else --
and in the latter case, what and why. Note that the frontend is delivered as a
\textbf{\texttt{.deb} installer} and collects the API URL and key in a
\textbf{connection dialog} at startup.

\subsection{Operation with Docker Compose}

List the \textbf{services} that will run in your \texttt{docker-compose.yml} (at
least \texttt{postgres} and \texttt{api}), where the credentials live
(\texttt{.env}, L 09) and how the backend is \textbf{deployed on the lecture
server}.

\subsection{Scope, milestones and risks}

Estimate the scope honestly: how many tables and endpoints? Name \textbf{2--3
milestones} (e.g.\ ``schema + seed data ready'', ``API reading/writing'',
``frontend operates the core endpoints'', ``documentation + video'') and
\textbf{one or two risks} that could make it fail. Keep the total within the
\textbf{$\approx$40\,h / 2-month} budget.

% ============================================================
\section{Formal requirements for the submission}
% ============================================================

\begin{itemize}[leftmargin=1.4em]
  \item \textbf{Format:} one PDF, 2--4 pages, with the sections from
        Section~5 (start from \texttt{proposal-template/}).
  \item \textbf{Submission:} by email to
        \href{mailto:sboekelmann@ep1.rub.de}{sboekelmann@ep1.rub.de}. Also place
        the PDF in your project repository (e.g.\ as \texttt{docs/proposal.pdf})
        once you have created one.
  \item \textbf{Subject / file name:} \texttt{term-project-proposal
        <surname>}.
  \item \textbf{Language:} English.
\end{itemize}

\begin{warnbox}
\textbf{Do not start the implementation} before the proposal is approved. Wait
for the lecturer's feedback -- it may still shift scope or data model and thus
save you unnecessary work.
\end{warnbox}

% ============================================================
\section{Acceptance criteria -- checklist}
% ============================================================

Your proposal is accepted if it satisfies the following points. Use the list as
a self-check before submitting.

\begin{itemize}[leftmargin=1.8em, label=$\square$]
  \item Application domain and benefit are described clearly.
  \item An ER sketch with cardinalities is present; at least one \texttt{N:M}
        relationship is included.
  \item A relational schema with primary and foreign keys is derived
        (3 to 6 tables, at least 3NF).
  \item At least one reading and one writing API endpoint are planned; JOIN and
        aggregation are located; write endpoints require the \texttt{X-API-Key}.
  \item The frontend and the endpoints it uses are sketched; it is delivered as
        a \texttt{.deb} with a connection dialog.
  \item The Docker Compose backend (postgres + api, \texttt{.env}) and its
        deployment on the lecture server are named.
  \item Scope, milestones and risks are realistic and fit the
        $\approx$40\,h / 2-month budget.
  \item The mandatory stack (PostgreSQL, FastAPI, Docker Compose, frontend) is
        fully covered.
\end{itemize}

% ============================================================
\section{A worked mini-example}
% ============================================================

So you can gauge the expected level, here is a strongly shortened example for
the domain \emph{``Club Library''}. Your own proposal is more detailed; this
only shows the direction. A full example documentation lives in
\texttt{example-documentation/}.

\begin{beispielbox}
\textbf{Domain.} A small club library manages books, members and loans. Members
want to see which books are available, and the librarian checks loans out and
back in.

\medskip
\textbf{Data model (excerpt).}
\begin{itemize}[leftmargin=1.4em, topsep=2pt]
  \item \texttt{member}(\underline{id}, name, email)
  \item \texttt{book}(\underline{id}, title, author, isbn, stock)
  \item \texttt{loan}(\underline{id}, \emph{member\_id}, \emph{book\_id},
        from, to, returned) -- realises the \texttt{N:M} relationship between
        member and book.
\end{itemize}

\textbf{API (excerpt).}
\begin{itemize}[leftmargin=1.4em, topsep=2pt]
  \item \texttt{GET /books} -- all books with the number of current loans
        (\emph{JOIN + aggregation}).
  \item \texttt{GET /members/\{id\}/loans} -- open loans of a member
        (\emph{JOIN}).
  \item \texttt{POST /loans} -- lend a book (\emph{requires X-API-Key}).
  \item \texttt{POST /returns} -- return a book (\emph{requires X-API-Key}).
\end{itemize}

\textbf{Frontend.} tkinter window with two tabs: ``Books'' (table + search
field) and ``Loans'' (form for lending/returning). Packaged as a \texttt{.deb};
asks for API URL and key at startup.

\medskip
\textbf{Operation.} \texttt{docker-compose.yml} with the services
\texttt{postgres} and \texttt{api}; passwords in \texttt{.env}; schema and seed
data via an init script; deployed on the lecture server.

\medskip
\textbf{Milestones.} (1)~schema + seed data; (2)~API reading/writing;
(3)~frontend operates both tabs, packaged as \texttt{.deb};
(4)~documentation + video. \textbf{Risk:} correct return handling when several
open loans of the same book exist at once.
\end{beispielbox}

% ============================================================
\section{Next steps}
% ============================================================

\begin{enumerate}[leftmargin=1.6em]
  \item Choose a domain that interests you -- you will work on it for up to two
        months.
  \item Sketch the ER model on paper and derive the schema.
  \item Copy \texttt{proposal-template/} and fill in the sections from
        Section~5, then check the checklist in Section~7.
  \item Send the PDF to the lecturer and wait for approval.
\end{enumerate}

\vspace{1em}
\begin{infobox}
\textbf{Further reading:} The exercises DBMS\_07 (PostgreSQL), DBMS\_08
(FastAPI), DBMS\_09 (Docker Compose) and lecture~10 (frontend) are your most
important references for the implementation. You need not learn anything new for
the proposal -- only apply the familiar to your own idea.
\end{infobox}

\end{document}
